```latex
\documentclass[11pt]{article}
\usepackage{amsmath,amssymb,amsthm,mathrsfs}

\usepackage[margin=1.2in]{geometry}

\newtheorem{theorem}{Theorem}
\newtheorem{lemma}{Lemma}
\newtheorem{proposition}{Proposition}
\newtheorem{corollary}{Corollary}
\theoremstyle{definition}
\newtheorem{definition}{Definition}
\theoremstyle{remark}
\newtheorem{remark}{Remark}

\newcommand{\D}{\mathcal{D}}
\newcommand{\Dist}{\mathcal{D}'}
\newcommand{\Lie}{\mathcal{L}}
\newcommand{\supp}{\operatorname{supp}}
\newcommand{\htop}{h_{\mathrm{top}}}

\begin{document}

\title{Globally hypoelliptic vector fields have zero topological entropy}
\author{}
\date{}
\maketitle

\begin{abstract}
Let $M$ be a compact connected smooth manifold and $X$ a smooth vector field with flow $(g_t)_{t\in\mathbb{R}}$.
Assume that the Lie derivative $\Lie_X$ is \emph{globally hypoelliptic} in the sense that whenever $u\in\Dist(M)$ and
$\Lie_X u=f\in C^\infty(M)$, then $u$ is a smooth function. We prove that the flow $g_t$ has zero topological entropy.
The proof uses (i) the fact that global hypoellipticity forces all $X$--invariant distributions to be smooth,
and (ii) Katok's horseshoe theorem, which yields periodic orbits under positive entropy for $C^{1+\alpha}$ diffeomorphisms.
\end{abstract}

\section{Preliminaries}

Throughout, $M$ denotes a compact connected $C^\infty$ manifold of dimension $d:=\dim M$, and
$X\in \mathfrak{X}^\infty(M)$ is a smooth vector field with (complete) flow $(g_t)_{t\in\mathbb{R}}$.

\subsection{Distributions and Lie derivative}

Let $\Dist(M)$ denote the space of distributions on $M$, i.e.\ continuous linear functionals on $\D(M)=C^\infty(M)$
(with the standard Fr\'echet topology); since $M$ is compact, the distinction between $C_c^\infty(M)$ and $C^\infty(M)$ is immaterial.

For $\varphi\in C^\infty(M)$ we write $X\varphi:=d\varphi(X)$.
For $u\in\Dist(M)$, define $\Lie_X u\in\Dist(M)$ by duality:
\begin{equation}\label{eq:LieDist}
  \langle \Lie_X u,\varphi\rangle := - \langle u, X\varphi\rangle,\qquad \forall \varphi\in C^\infty(M).
\end{equation}
This agrees with the usual Lie derivative on smooth functions (viewed as distributions) by integration by parts
after choosing any smooth volume density.

\begin{definition}[Global hypoellipticity for $\Lie_X$]\label{def:GH}
We say that $\Lie_X$ is \emph{globally hypoelliptic} if for every $u\in \Dist(M)$,
\[
\Lie_X u\in C^\infty(M)\quad \Longrightarrow\quad u\in C^\infty(M).
\]
Equivalently (and matching the hypothesis of the problem), if $f\in C^\infty(M)$ and $u\in\Dist(M)$ satisfy
$\Lie_X u=f$, then $u$ is a smooth function.
\end{definition}

\subsection{Topological entropy for flows}

For a continuous flow $(g_t)$, one defines $\htop(g_t)$ in several equivalent ways. One standard convention is
\[
\htop((g_t)_{t\in\mathbb{R}}):=\htop(g_1),
\]
the topological entropy of the time--one map. Another standard fact is linear scaling in time:
\[
\htop(g_t)=|t|\htop(g_1),\qquad t\in\mathbb{R},
\]
see e.g.\ \cite[Theorem 7.2]{Walters} or Bowen's work on entropy for flows.
In particular, $\htop(g_t)=0$ for all $t$ if and only if $\htop(g_1)=0$.

\section{Consequences of global hypoellipticity}

\begin{lemma}\label{lem:kernelSmooth}
Assume $\Lie_X$ is globally hypoelliptic. Then every distributional solution of $\Lie_X u=0$ is smooth. In particular,
every $X$--invariant distribution is a smooth function (as a distribution).
\end{lemma}

\begin{proof}
Take $f\equiv 0\in C^\infty(M)$. The hypothesis gives:
if $u\in\Dist(M)$ and $\Lie_X u=0$, then $u\in C^\infty(M)$.

If $u$ is $X$--invariant in the sense $(g_t)_*u=u$ for all $t$, then differentiating at $t=0$ in the distributional sense
gives $\Lie_X u=0$; equivalently, by \eqref{eq:LieDist}, invariance implies $\langle u, X\varphi\rangle=0$ for all $\varphi$.
Thus $u$ is smooth.
\end{proof}

\begin{lemma}\label{lem:noZeros}
Assume $\Lie_X$ is globally hypoelliptic. Then $X$ has no zeroes.
\end{lemma}

\begin{proof}
Suppose $X(p)=0$ for some $p\in M$.
Let $\delta_p\in\Dist(M)$ be the Dirac distribution at $p$: $\langle \delta_p,\varphi\rangle=\varphi(p)$.
Then for every $\varphi\in C^\infty(M)$,
\[
\langle \Lie_X \delta_p,\varphi\rangle = -\langle \delta_p, X\varphi\rangle = -(X\varphi)(p)=0
\]
since $X(p)=0$. Hence $\Lie_X\delta_p=0$ (a smooth function).
By global hypoellipticity, $\delta_p$ must be a smooth function distribution, which is impossible.
Thus $X$ is nowhere vanishing.
\end{proof}

\begin{lemma}\label{lem:orbitMeasureInvariant}
Let $\gamma$ be a periodic orbit of the flow $g_t$ with period $T>0$. Define a distribution $\mu_\gamma$ by
\begin{equation}\label{eq:orbitmeasure}
 \langle \mu_\gamma,\varphi\rangle := \frac{1}{T}\int_0^T \varphi(g_t(x))\,dt,\qquad \varphi\in C^\infty(M),
\end{equation}
where $x\in \gamma$. Then $\mu_\gamma$ is $g_t$--invariant and satisfies $\Lie_X \mu_\gamma=0$.
\end{lemma}

\begin{proof}
The definition is independent of the choice of $x\in\gamma$ because the parametrization is periodic with period $T$.
Invariance under the flow is immediate:
\[
\langle (g_s)_*\mu_\gamma,\varphi\rangle
=\langle \mu_\gamma,\varphi\circ g_s\rangle
=\frac{1}{T}\int_0^T \varphi(g_s(g_t(x)))\,dt
=\frac{1}{T}\int_0^T \varphi(g_{t+s}(x))\,dt
=\langle \mu_\gamma,\varphi\rangle
\]
after changing variables $t\mapsto t+s$ and using periodicity.
Differentiating invariance at $s=0$ yields $\Lie_X\mu_\gamma=0$; equivalently, compute directly:
\[
\langle \Lie_X\mu_\gamma,\varphi\rangle
= -\langle \mu_\gamma, X\varphi\rangle
= -\frac{1}{T}\int_0^T (X\varphi)(g_t(x))\,dt
= -\frac{1}{T}\int_0^T \frac{d}{dt}(\varphi(g_t(x)))\,dt
=0.
\]
\end{proof}

\begin{lemma}\label{lem:smoothDistributionSupport}
Let $u\in C^\infty(M)$ be viewed as a distribution. If $u\not\equiv 0$, then its support $\supp(u)$ has nonempty interior.
Consequently, any nonzero distribution whose support has empty interior cannot be a smooth function distribution.
\end{lemma}

\begin{proof}
If $u\not\equiv 0$, there exists $p\in M$ with $u(p)\neq 0$.
By continuity, there exists an open neighborhood $U$ of $p$ such that $u|_U\neq 0$.
Thus $\{u\neq 0\}$ is a nonempty open set. The support $\supp(u)$ is the closure of $\{u\neq 0\}$, hence contains $U$,
and therefore has nonempty interior.
\end{proof}

\begin{lemma}\label{lem:periodicOrbitEmptyInterior}
If $d=\dim M\ge 2$ and $\gamma\subset M$ is a periodic orbit of a nowhere--vanishing vector field,
then $\gamma$ has empty interior in $M$.
\end{lemma}

\begin{proof}
Since $X$ is nowhere vanishing (Lemma \ref{lem:noZeros}), every orbit is a one--dimensional immersed submanifold;
a periodic orbit is a compact one--dimensional embedded submanifold diffeomorphic to $S^1$.
In a coordinate chart $(U,\psi)$ with $\psi:U\to \mathbb{R}^d$, the set $\psi(\gamma\cap U)$ is contained in a
$1$--dimensional embedded $C^\infty$ submanifold of $\mathbb{R}^d$. Such a set cannot contain any nonempty open subset of
$\mathbb{R}^d$ when $d\ge 2$. Hence $\gamma$ cannot contain any nonempty open subset of $M$, i.e.\ $\gamma$ has empty interior.
\end{proof}

\section{Positive entropy forces periodic orbits (Katok)}

We use a deep theorem due to Katok which, combined with the variational principle, implies that positive \emph{topological}
entropy for a $C^{1+\alpha}$ diffeomorphism forces periodic points.

\begin{theorem}[Variational principle]\label{thm:variational}
Let $f$ be a continuous map on a compact metric space. Then
\[
\htop(f)=\sup_{\mu} h_\mu(f),
\]
where the supremum is over all $f$--invariant Borel probability measures $\mu$, and $h_\mu(f)$ is the measure--theoretic entropy.
\end{theorem}

A proof can be found in \cite[Theorem 8.6]{Walters}.

\begin{theorem}[Katok's horseshoe theorem]\label{thm:katok}
Let $f$ be a $C^{1+\alpha}$ diffeomorphism of a compact smooth manifold ($\alpha>0$).
If $\mu$ is an ergodic $f$--invariant probability measure with $h_\mu(f)>0$, then for every $\varepsilon>0$
there exists a compact $f$--invariant hyperbolic set $\Lambda$ (a horseshoe) such that
$\htop(f|_\Lambda)\ge h_\mu(f)-\varepsilon$.
In particular, $f$ has a periodic point.
\end{theorem}

This is proved in \cite[Theorem~1]{KatokIHES} (see also \cite[Chapter~18]{KH} for modern expositions).

\begin{remark}\label{rem:periodicPoint}
The ``in particular'' follows because every horseshoe contains periodic points (indeed, it is semi-conjugate to a
subshift of finite type, which has periodic points). Standard references include \cite[\S 18.6]{KH}.
\end{remark}

\section{Main result}

\begin{theorem}\label{thm:main}
Let $M$ be a compact connected smooth manifold and $X$ a smooth vector field with flow $(g_t)$.
Assume that $\Lie_X$ is globally hypoelliptic in the sense of Definition~\ref{def:GH}.
Then the flow $(g_t)$ has zero topological entropy, i.e.\ $\htop(g_1)=0$ and hence $\htop(g_t)=0$ for all $t\in\mathbb{R}$.
\end{theorem}

\begin{proof}
We prove $\htop(g_1)=0$. Suppose, for contradiction, that
\begin{equation}\label{eq:positive}
 \htop(g_1)>0.
\end{equation}

\smallskip\noindent\textbf{Step 1: Reduce to the case $\dim M\ge 2$.}
If $\dim M=1$, then $M$ is a disjoint union of circles; since $M$ is connected, $M\simeq S^1$.
The time--one map $g_1$ is a $C^\infty$ diffeomorphism of the circle. Every circle homeomorphism (hence every circle diffeomorphism)
has zero topological entropy; see \cite[Example 7.3.3 and Theorem 7.4]{Walters}.
Thus $\htop(g_1)=0$, contradicting \eqref{eq:positive}. Therefore we may assume $\dim M=d\ge 2$.

\smallskip\noindent\textbf{Step 2: Positive entropy gives a periodic orbit.}
By the variational principle (Theorem~\ref{thm:variational}) applied to $g_1$, there exists a $g_1$--invariant ergodic probability
measure $\mu$ such that $h_\mu(g_1)>0$ (since $\htop(g_1)>0$).
The diffeomorphism $g_1$ is $C^\infty$, hence $C^{1+\alpha}$ for some $\alpha>0$, so Katok's theorem (Theorem~\ref{thm:katok})
implies that $g_1$ has a periodic point $p\in M$:
there exists $n\ge 1$ such that
\[
g_1^n(p)=p.
\]
But $g_1^n=g_n$ (flow property), so $g_n(p)=p$, hence $p$ lies on a periodic orbit $\gamma=\{g_t(p):t\in\mathbb{R}\}$ of the flow.

\smallskip\noindent\textbf{Step 3: The periodic orbit yields a non-smooth invariant distribution.}
Let $T>0$ be the minimal period of $\gamma$.
Define the orbit measure $\mu_\gamma$ by \eqref{eq:orbitmeasure}.
By Lemma~\ref{lem:orbitMeasureInvariant}, $\Lie_X\mu_\gamma=0$.

By global hypoellipticity (with $f\equiv 0$), Lemma~\ref{lem:kernelSmooth} gives that $\mu_\gamma$ must be a smooth function
(as a distribution).

On the other hand, $\supp(\mu_\gamma)\subseteq \gamma$, hence $\supp(\mu_\gamma)\subseteq \gamma$ has empty interior in $M$
by Lemma~\ref{lem:periodicOrbitEmptyInterior} (recall $d\ge 2$).
But $\mu_\gamma\neq 0$, and if a distribution is represented by a nonzero smooth function, its support has nonempty interior
(Lemma~\ref{lem:smoothDistributionSupport}).
This contradiction shows that \eqref{eq:positive} is impossible.

Therefore $\htop(g_1)=0$.

\smallskip\noindent\textbf{Step 4: Entropy of the flow.}
By the standard scaling property for entropy of flows (see \cite[Theorem 7.2]{Walters}),
$\htop(g_t)=|t|\htop(g_1)=0$ for all $t\in\mathbb{R}$.
\end{proof}

\section*{Acknowledgments}
The argument is standard: global hypoellipticity excludes singular invariant distributions, while positive entropy forces periodic
orbits and hence singular invariant measures.

\begin{thebibliography}{99}

\bibitem{KatokIHES}
A.~Katok,
\newblock \emph{Lyapunov exponents, entropy and periodic orbits for diffeomorphisms}.
\newblock Publications Math\'ematiques de l'IH\'ES \textbf{51} (1980), 137--173.

\bibitem{KH}
A.~Katok and B.~Hasselblatt,
\newblock \emph{Introduction to the Modern Theory of Dynamical Systems}.
\newblock Encyclopedia of Mathematics and its Applications, vol.~54, Cambridge University Press, 1995.

\bibitem{Walters}
P.~Walters,
\newblock \emph{An Introduction to Ergodic Theory}.
\newblock Graduate Texts in Mathematics, vol.~79, Springer, 1982.

\end{thebibliography}

\end{document}
```